%----------------------------------------------------------------------------------------
% PACKAGES AND DOCUMENT CONFIGURATIONS
%----------------------------------------------------------------------------------------

\documentclass{article}
%packages to use various commands. You may need more packages depending on the commands you'd like to use, but it's unlikely.
\usepackage{fullpage}
\usepackage{graphicx}
\usepackage{fancyhdr}
\usepackage{hyperref}
\usepackage{amsmath}
\usepackage{tikz}
\usepackage{titling}
\pagestyle{fancy}

\usepackage{listings}
\usepackage{xcolor} % For colors
\usepackage{amssymb}
\usepackage{booktabs}
\usepackage{siunitx}
\usepackage[version=4]{mhchem}

% Define a custom style
\lstset{
    basicstyle=\ttfamily\small,
    keywordstyle=\color{blue},
    commentstyle=\color{green!60!black},
    stringstyle=\color{red},
    breaklines=true,
    frame=single,
    numbers=left,
    numberstyle=\tiny\color{gray},
    captionpos=b % 'b' for bottom, 't' for top
}



\setlength{\droptitle}{-10em} 

\renewcommand{\headheight}{32px}
\renewcommand{\headrulewidth}{0.75pt}
\setlength{\headsep}{20px}
\setlength{\parskip}{8px}
\setlength{\topmargin}{-18px}
\fancyhead[LO]{\small McMaster University \\ Math 3MB3}
\fancyhead[RO]{\small Final Project}

% a new command that makes a better +/- sign than LaTeX provides
\newcommand{\ppm}{\mathbin{\tikz [x=1.4ex,y=1.4ex,line width=.1ex] \draw (0.0,0) -- (1.0,0) (0.5,0.08) -- (0.5,0.92) (0.0,0.5) -- (1.0,0.5);}}


%----------------------------------------------------------------------------------------
% DOCUMENT INFORMATION
%----------------------------------------------------------------------------------------

\title {Final Project \\  Math 3MB3 } % Title

\author{ J. \textsc{McElheron et al.}} % Author name

\date{\today} % Date for the report

\begin{document}

\begin{titlepage}
\maketitle % Insert the title, author and date
\thispagestyle{empty}
\end{titlepage}

%----------------------------------------------------------------------------------------
% SECTION 1
%----------------------------------------------------------------------------------------
\section{Section 1}
%----------------------------------------------------------------------------------------
% SECTION 2
%----------------------------------------------------------------------------------------
\section{Section 2}

%----------------------------------------------------------------------------------------
% SECTION 3
%----------------------------------------------------------------------------------------
\section{Section 3}

%-------------------------------------------------------------------------
%References
%-------------------------------------------------------------------------
\newpage
\section{References}

\end{document}

